%%%%%%%%%%%%%%%%%%%%%%%%%%%%%%%%%%%%%%%%%%%%%%%%%%%%%%%%%%%%%%%%%%%%%%%%%%%%%%%%
\section*{Pair Programming}
{   
    \usebackgroundtemplate{
        \vbox to \paperheight{\vfil\hbox to \paperwidth{\hfil\includegraphics[width=\paperwidth]{humor/pair_programming.png}\hfil}\vfil}
        % source: https://en.m.wikipedia.org/wiki/File:Text-x-python.svg
    }
    \frame{
        \frametitle{What is Pair Programming?}
        \vspace{8.5em}
        \begin{mdframed}[tikzsetting={draw=white,fill=white,fill opacity=0.8,
                line width=0pt},backgroundcolor=none,leftmargin=0,
            rightmargin=50,innertopmargin=4pt,roundcorner=10pt]
            \bcquestion What is Pair Programming about?
                \begin{itemize}
                    \item you take turns being the \textbf{Driver} (controls keyboard and mouse) and the \textbf{Navigator} (observe, catch bugs early, suggest next steps)
                    \item you swap after each exercise or after 15 minutes -- whichever comes first
                    \item the idea is to get you talking and help each other
                \end{itemize}

        \end{mdframed}
        \vspace{10mm}
        \hfill{\tiny sketch by C. Meesters, coal and sanguine, CC-BY-4.0}
    }
}